\section{联合、边缘与条件分布}
\label{sec:05-Probability/05-Joint-Distributions/01}

你手上有两份数据：城市每天是否下雨的记录，以及每天道路是否拥堵的记录。单看任何一份都能算出一个概率——\(P(\text{下雨})=0.3\)，\(P(\text{拥堵})=0.4\)。然后老板问你一个听起来很简单的问题：``下雨天的拥堵概率是多少？''

你盯着这两个数字。你只有它们各自的频率，不知道它们是怎么配对出现的。30\% 的下雨天里，拥堵发生在其中的多少天？可能是全部，也可能只有一小部分。两个边缘概率无法回答这个问题。

这不是数据不够的问题——这是你还没把数据放在一起看的问题。单个变量的概率分布告诉你``X 取某个值的频率''，但 X 和 Y 取值的组合频率，需要另一种结构来记录。

\subsection{联合分布保留配对信息}

好，假设我们把每天的数据配对记录。某天既下雨又拥堵，某天只下雨不拥堵，某天只拥堵不下雨，某天两者皆无。四种可能的世界。

离散情形下，联合概率质量函数（joint PMF）把每一个 \((x,y)\) 组合映射到一个概率：

\[
p_{X,Y}(x,y)=P(X=x,Y=y),
\]

并且所有组合的概率加起来必须等于 1：

\[
p_{X,Y}(x,y)\ge0,\qquad
\sum_x\sum_y p_{X,Y}(x,y)=1.
\]

用一张 \(2\times2\) 表格来存放会很直观：

\[
\begin{array}{c|cc|c}
&Y=0\ (\text{不堵})&Y=1\ (\text{堵})&\text{行和}\\
\hline
X=0\ (\text{晴})&0.4&0.1&0.5\\
X=1\ (\text{雨})&0.2&0.3&0.5\\
\hline
\text{列和}&0.6&0.4&1
\end{array}
\]

表格里的四个数字就是联合分布。每一格是一类``配对事件''。晴且不堵占 40\%，雨且堵占 30\%，等等。

行和与列和就是两个变量各自的分布——你看，把联合表格沿行方向加总，得到的是 X 自己的分布，这跟你一开始手里只有两个单独的频率，信息量完全不同。单独频率可以对应无数种内部填法；联合表格把内部填法唯一确定了下来。

连续变量的情况是类似的想法，只不过把求和换成积分。如果 \((X,Y)\) 存在联合密度 \(f_{X,Y}\)，那么 \((X,Y)\) 落在平面区域 \(A\subseteq\R^2\) 里的概率是

\[
P\bigl((X,Y)\in A\bigr)
=\iint_A f_{X,Y}(x,y)\,dx\,dy.
\]

密度非负，且在整个平面上的积分等于 1。

这里有一个容易被跳过的细节：支撑区域。如果联合密度只在三角形 \(0<y<x<1\) 上非零，那么积分上下限必须刻画这个三角形，不能想当然地写成从 0 到 1 的双重积分然后塞进一个常数。三角区域的二重积分上下限是 \(\int_{x=0}^1\int_{y=0}^x\)，或者等价地 \(\int_{y=0}^1\int_{x=y}^1\)。如果写错积分的次序和范围，后面的所有推导全废。这是联合分布区别于两个独立单变量分布的第一道关卡：支撑形状不一定是矩形。

\subsection{边缘化是把另一个变量忘掉}

联合分布里有两个变量。有时你只关心其中一个，另一个不感兴趣。那你需要把不感兴趣的变量``消掉''。

离散时：

\[
p_X(x)=\sum_y p_{X,Y}(x,y).
\]

连续时：

\[
f_X(x)=\int_{-\infty}^{\infty}f_{X,Y}(x,y)\,dy.
\]

这个操作叫边缘化（marginalization）。名字的来源就是那张表格：你沿着一行把概率加总，写在右边空白处（margin），得到的就是 X 的边缘分布。

在上面的表格里，

\[
P(X=1)=0.2+0.3=0.5.
\]

雷雨天的总概率 50\%，不论是否拥堵。

但这里有一个很重要的认知节点：一旦你只保留行和与列和，格子内部怎么分配的信息就永久丢失了。无数张不同的内部填充表格可以拥有完全相同的边缘分布。边缘化是不可逆操作。你之前遇到的困境——从 \(P(\text{下雨})\) 和 \(P(\text{拥堵})\) 推不出 \(P(\text{下雨且拥堵})\)——正是因为边缘化丢掉了配对信息。

\subsection{条件分布是归一化的切片}

有了联合表格，就可以回答老板的问题了。``已知下雨，拥堵的概率''意味着：我们只看 X=1 那一行，把这一行当作新的概率空间，重新归一化。

离散情形，只要 \(p_Y(y)>0\)：

\[
p_{X\mid Y}(x\mid y)
=\frac{p_{X,Y}(x,y)}{p_Y(y)}.
\]

在表格中计算给你看：

\[
P(X=1 \mid Y=1)=\frac{P(X=1,Y=1)}{P(Y=1)}=\frac{0.3}{0.4}=0.75.
\]

固定 \(Y=1\)，取出表格的第二列：\((0.1, 0.3)\)。这一列的和是 0.4。除以 0.4 后变成 \((0.25, 0.75)\)，总和回到 1。已知拥堵的条件下，下雨的概率是 75\%，远高于不下雨时的 25\%。而边缘概率 \(P(X=1)=0.5\)——条件分布给出了完全不同的图景。

这个操作可以反过来做。如果你的问题是``已知下雨，拥堵概率''，就固定行、归一化：

\[
P(Y=1 \mid X=1)=\frac{0.3}{0.5}=0.6.
\]

两种条件分布通常不相等。条件方向很关键。

连续情形也是类似的比例结构，在 \(f_Y(y)>0\) 处：

\[
f_{X\mid Y}(x\mid y)
=\frac{f_{X,Y}(x,y)}{f_Y(y)}.
\]

你可能会觉得这里有点微妙：连续变量取单点的概率是 0，\(P(Y=y)=0\)，那``给定 \(Y=y\)''算什么？条件密度的定义绕过了这个问题：它不是直接用条件概率公式，而是通过联合密度与边缘密度的比值来定义。可以把它理解为越来越窄的 Y 邻域上条件概率的极限，不过严格的存在性需要正则条件概率的框架来保证。对大部分应用来说，直接按比值计算即可。

\subsection{一个三角形支撑的例子}

理论讲到这里，如果不亲手算一个连续例子，很难真正消化积分上下限的问题。

设 \((X,Y)\) 在三角形区域

\[
0<y<x<1
\]

上具有常数密度 \(c\)，其余地方密度为零。三角形面积是 \(1/2\)，所以

\[
c\cdot\frac12=1,\qquad c=2.
\]

现在来求 X 的边缘密度。固定一个 \(x\in(0,1)\)，在三角形内部，y 的取值范围是从 0 到 \(x\)。把 y 积掉：

\[
f_X(x)=\int_0^x2\,dy=2x,\qquad 0<x<1.
\]

反过来，固定一个 \(y\in(0,1)\)，x 的取值范围是从 y 到 1：

\[
f_Y(y)=\int_y^1 2\,dx=2(1-y),\qquad 0<y<1.
\]

注意 \(f_X\) 和 \(f_Y\) 的形状完全不同：一个是随 x 增大而线性增长的，另一个是随 y 增大而线性衰减的。虽然联合密度在三角形上处处相等，但边缘化之后的不对称来自三角形边界的不对称——x 方向的``厚度''随 x 增大而增大，y 方向的``厚度''随 y 增大而减小。

条件密度：

\[
f_{Y\mid X}(y\mid x)
=\frac{2}{2x}
=\frac1x,\qquad 0<y<x.
\]

给定 \(X=x\) 后，Y 的分布在 \((0,x)\) 上是均匀的。这很合理：联合密度在三角形内是常数，所以固定 x 后，允许的 y 区间上密度也是常数，归一化后就得到均匀分布。

这个例子的要点不是计算本身，而是一个操作纪律：画支撑区域，在图上标出积分变量的上下限，然后再写积分。如果不画图，凭感觉写 \(\int_0^1\int_0^1\)，整个过程从一开始就是错的。

\subsection{联合 CDF 的角色}

联合累积分布函数

\[
F_{X,Y}(x,y)=P(X\le x,Y\le y)
\]

是一个更底层的定义工具。它的好处是：不管变量是离散、连续、还是混合类型，不管密度是否存在，联合 CDF 都能定义，并能唯一确定联合分布。

但实际计算中，离散表格和联合密度通常比 CDF 更直观。CDF 的角色更像是``法律上的最终依据''——当密度不存在或公式打架时，回到 CDF。

下一节会追问一个很自然的问题：什么情况下，联合分布恰好等于边缘分布的乘积？那也就是在问，两个随机变量什么时候``没有关系''——这正是独立性的精确定义。
